\chapter{III. PHẦN KẾT LUẬN}

\subsection{1. Kết quả đạt được}

\indent\indent Đề tài \textbf{``Phát hiện và phân loại đa lớp chữ bị làm giả trong ảnh tài liệu sử dụng mô hình học sâu đa nhiệm EfficientNet-B0''} đã hoàn thành toàn diện các mục tiêu đề ra với những đóng góp cụ thể:

\begin{itemize}[label={--}, itemsep=3pt]
    \item \textbf{Đề xuất thành công kiến trúc đa nhiệm gọn nhẹ LightweightDTD:} Xây dựng mô hình với tổng số chỉ \textbf{4,155,796 tham số (~4.15M)}, kết hợp hài hòa giữa nhánh không gian EfficientNet-B0 và nhánh tần số biến đổi trực chuẩn BlockDCT $8 \times 8$.
    \item \textbf{Thiết kế bộ giải mã LiteMID và khối chú ý scSE tối ưu:} Ứng dụng tích chập Depthwise Separable Convolutions để lồng ghép hiệu quả các biểu diễn đa tỉ lệ từ $H/4$ đến $H/32$, giúp tinh lọc mặt nạ phân đoạn với chi phí tính toán thấp.
    \item \textbf{Cơ chế học đa nhiệm toàn diện:} Mô hình vừa phân đoạn chính xác vùng điểm ảnh bị làm giả, vừa phân loại đồng thời 6 đặc trưng pháp y vi mô (\textit{edge\_artifacts}, \textit{unnatural\_texture}, \textit{inconsistent\_alignment}, \textit{inconsistent\_font}, \textit{lack\_of\_integration}, \textit{text\_incoherence}).
    \item \textbf{Đạt các chỉ số thực nghiệm xuất sắc trên Minh-TFR TestSet:}
    \begin{itemize}[label={+}, itemsep=1pt]
        \item Độ chính xác điểm ảnh (\textbf{Pixel Accuracy}): đạt \textbf{99.68\%}.
        \item Tỷ lệ bắt trúng nhãn (\textbf{Top-3 Class Hit Rate}): đạt \textbf{98.30\%}.
        \item Độ chuẩn xác (\textbf{Precision Micro}): đạt \textbf{85.25\%}.
        \item Độ trùng khớp (\textbf{IoU Micro}): đạt \textbf{58.39\%}.
    \end{itemize}
\end{itemize}

\subsection{2. Hạn chế và Hướng phát triển}

\subsubsection{2.1. Hạn chế của đề tài}
\begin{itemize}[label={--}, itemsep=2pt]
    \item \textbf{Độ phân giải cố định:} Mô hình hiện tại huấn luyện trên các ảnh kích thước chuẩn $224 \times 224$ pixel. Khi áp dụng cho các trang tài liệu khổ lớn A4 độ phân giải cao ($2000 \times 3000$ pixel), cần phải cắt nhỏ thành các khối (patches) để xử lý.
    \item \textbf{Độ nhạy với ảnh nén mất mát cực độ:} Trong các trường hợp tài liệu chụp qua camera điện thoại bị mờ rung mạnh hoặc chuyển tiếp nhiều lần qua các mạng xã hội với tỉ lệ nén rất cao ($Q < 30$), tín hiệu tần số DCT có thể bị suy biến, ảnh hưởng đến độ nhạy phát hiện.
\end{itemize}

\subsubsection{2.2. Hướng nghiên cứu và phát triển tiếp theo}
\begin{itemize}[label={--}, itemsep=2pt]
    \item \textbf{Mở rộng cơ chế End-to-End toàn trang tài liệu:} Phát triển mô-đun phát hiện vùng văn bản nghi vấn toàn trang dựa trên mạng trượt cửa sổ thích ứng (Sliding Window / Patch-based aggregation) để tự động hóa hoàn toàn quy trình xử lý.
    \item \textbf{Tích hợp suy luận biên trên thiết bị di động:} Ứng dụng các kỹ thuật lượng tử hóa (INT8 Quantization, ONNX Runtime, TensorRT) để triển khai trực tiếp mô hình vào các ứng dụng máy quét trên điện thoại thông minh phục vụ người dùng cuối.
    \item \textbf{Đa dạng hóa tập dữ liệu văn bản tiếng Việt:} Thu thập và mở rộng tập dữ liệu cho các tài liệu hành chính tại Việt Nam như căn cước công dân gắn chip, sao kê ngân hàng và hóa đơn giá trị gia tăng điện tử.
\end{itemize}
